\section{Evaluation}
\label{sec:evaluation}

\subsection{Protocol}

The still-image check uses the Human Action Recognition set released by
Bingsu. The published test split has placeholder labels, so the sample is
drawn from the labeled train split: 15 classes, 4 images per class, seed 0,
60 images. The models are not trained on it. Each model sees one image and
writes a free phrase with the same short-action prompt. The predicted class
is the class sentence with the highest text-encoder cosine. An earlier
matcher based on character overlap, with a bonus when the class name
appeared in the phrase, is kept only as a reference.

This protocol scores \eqref{eq:class}. It does not score
\eqref{eq:keep}, \eqref{eq:mad}, or \eqref{eq:fuse}. There is no ongoing
action across the 60 images, and Qwen sees one frame rather than three.
A phrase can describe the picture in ordinary English and still miss the
class, because the class sentence is a single canonical line. ``Typing''
sits near both texting and using a laptop. The cosine can prefer either
neighbor.

\subsection{Phrase-to-class accuracy}

\begin{table}[t]
\caption{Phrase-to-class accuracy on 60 still images.}
\label{tab:still}
\centering
\begin{tabular}{lcc}
\toprule
Matcher & SmolVLM & Qwen \\
\midrule
Character overlap & 27/60 & 30/60 \\
Text-encoder cosine & 38/60 & 40/60 \\
\bottomrule
\end{tabular}
\end{table}

Table~\ref{tab:still} shows that the text encoder recovers paraphrases that
character overlap misses. ``Riding bike'' matches cycling, and ``Drinking
wine'' matches drinking, without the class name having to appear as a
substring. The gain is 11 images for SmolVLM and 10 images for Qwen. Qwen
is ahead of SmolVLM by two images under the cosine matcher. On this sample
the larger model is not a large gain once both models are allowed to speak
in their own words and a text encoder absorbs the wording.

Table~\ref{tab:class} breaks the cosine score down by class. Both models
are correct on all four running images and fail most texting and
laptop images. SmolVLM is correct on all four calling images. Qwen writes
``talking'' or ``using a phone'' on several of those images, and the cosine
then prefers texting or calling inconsistently, so Qwen scores 1/4 on
calling. Qwen is correct on all four hugging, cycling, drinking, and
sleeping images, where SmolVLM drops two or one. The errors cluster on
actions that share a posture or an object: texting, using a laptop, and
sitting; smiling against laughing; talking against calling. Several misses
are confident. SmolVLM assigns 0.90 to ``Taking picture'' on a texting
image and 0.84 to ``Reading'' on a sleeping image. A high minimum token
probability means the generator was sure of its tokens. It does not mean
the phrase is the class sentence.

\begin{table}[t]
\caption{Cosine matches per class, out of four images.}
\label{tab:class}
\centering
\footnotesize
\begin{tabular}{lccc}
\toprule
Class & SmolVLM & Qwen & Mixed at 0.5 \\
\midrule
calling & 4 & 1 & 4 \\
clapping & 2 & 3 & 3 \\
cycling & 3 & 4 & 4 \\
dancing & 2 & 1 & 2 \\
drinking & 3 & 4 & 3 \\
eating & 4 & 3 & 4 \\
fighting & 2 & 3 & 3 \\
hugging & 2 & 4 & 3 \\
laughing & 3 & 3 & 2 \\
listening to music & 2 & 2 & 2 \\
running & 4 & 4 & 4 \\
sitting & 2 & 2 & 2 \\
sleeping & 3 & 4 & 3 \\
texting & 1 & 1 & 1 \\
using a laptop & 1 & 1 & 1 \\
\midrule
Total & 38 & 40 & 41 \\
\bottomrule
\end{tabular}
\end{table}

\subsection{Confidence as a router}

Deferring to Qwen when the SmolVLM confidence is below a threshold uses the
cosine predictions. Table~\ref{tab:route} reports three thresholds. At 0.5,
the mixed result is 41/60 and 26 of the 60 images take the Qwen phrase. At
0.3 the result is 42/60 with 12 images deferred. At 0.7 it is 43/60 with 41
images deferred. The gain over Qwen alone is at most three images. Raising
the threshold uploads more and barely moves the class score.

\begin{table}[t]
\caption{Defer to Qwen when SmolVLM confidence is below the threshold.}
\label{tab:route}
\centering
\begin{tabular}{lcc}
\toprule
Threshold & Correct & Images sent to Qwen \\
\midrule
0.3 & 42/60 & 12 \\
0.5 & 41/60 & 26 \\
0.7 & 43/60 & 41 \\
\bottomrule
\end{tabular}
\end{table}

The confidence values explain the flat curve. On the CUDA run, SmolVLM's
mean confidence is 0.577 on the 38 correct phrases and 0.484 on the 22
misses. At 0.5 the high-confidence slice is 22/34 correct and the
low-confidence slice is 16/26 correct. The two rates are close, so the
gate does not separate the phrases that need a second model from the
phrases that are already right. Qwen's own confidence separates a little
more, 25/34 above 0.5 against 15/26 below, but the router never reads
Qwen's confidence. It has to decide before the upload.

The two models are complementary on the class labels, and the gate does
not harvest that complementarity. Both phrases match the label on 30
images. SmolVLM alone is right on 8, Qwen alone is right on 10, and both
are wrong on 12. An oracle that kept the correct phrase whenever either
model produced one would score 48/60. The best threshold in
Table~\ref{tab:route} scores 43/60. Laughing is a concrete loss: each model
alone scores 3/4, and the mix at 0.5 scores 2/4, because the gate hands
Qwen an image that SmolVLM already had right. The live rule in
\eqref{eq:fuse} can also refuse a Qwen phrase that drifts away from the
ongoing action. The 60-image mix never applies that second test.

\subsection{Offline generation time}

Table~\ref{tab:time} is the wall time of \texttt{eval\_har.py} on CUDA
after the weights are loaded. SmolVLM is float32 in that harness. Qwen is
bfloat16, which is the server configuration. These numbers are not the
phone. The phone runs 4-bit SmolVLM in the browser, and on the development
phone that generation is under three seconds. The laptop figures say only
that, once the picture has arrived, one-image Qwen generation is usually
about 0.2~s, with one outlier at 1.2~s. A two-frame spot check outside this
table took 1.4~s. The confirmation delay in \eqref{eq:conf} is therefore
dominated by the phone generation period, the uplink, and any queue, not by
a multi-second laptop decode of one JPEG. That split is why the November
predictor has to be trained on traces that vary JPEG size and laptop load,
rather than on a single constant.

\begin{table}[t]
\caption{Offline CUDA generation time per still image, in milliseconds.}
\label{tab:time}
\centering
\begin{tabular}{lcccc}
\toprule
Model & Min & Median & Mean & Max \\
\midrule
SmolVLM, float32 & 468 & 538 & 582 & 1040 \\
Qwen3.5-2B, bfloat16 & 184 & 210 & 235 & 1184 \\
\bottomrule
\end{tabular}
\end{table}

\subsection{Three-frame replay}

The live rule is replayed on the UCF101 recognition test list and on the
HMDB51 test list packaged with the Sina272 release. UCF101 contributes
3783 clips and 101 classes. HMDB51 contributes 1530 clips, 51 classes, and
30 clips per class. Each clip is reduced to three JPEGs, longest side 512.
When the clip lasts at least 2.2 seconds the frames sit one second before
the middle, at the middle, and one second after. A shorter clip spreads
the three indexes across its length. The replay then walks those frames in
order with Algorithm~\ref{alg:live}. SmolVLM-only uses the quiet skip and
never uploads. Qwen-only sees all three frames in one call. The hybrid
branch is the live rule, with \(\tau_{\mathrm{u}}=0.5\) and
\(\tau_{\mathrm{s}}=0.70\).

The predicted class is again \eqref{eq:class}. The class sentence is the
template ``a person is'' followed by the class name, with camel case split
into words, so the sentence is not always grammatical. The models are not
trained on either set. Latency sums the pixel check, the SmolVLM calls, the
Qwen calls, and the phrase-encoder calls used on that clip. It does not
include Wi-Fi time. SmolVLM is float32 on CUDA. Qwen is bfloat16. These
times are not the phone. The three frames are taken around the middle of a
finished clip, so the replay does not measure an action's start or end,
and it is not the official accuracy of either dataset.

\begin{table}[t]
\caption{Three-frame replay. Uploads are Qwen calls per clip. Skipped frames bypass SmolVLM.}
\label{tab:video}
\centering
\footnotesize
\setlength{\tabcolsep}{4pt}
\begin{tabular}{llcccc}
\toprule
Set & Branch & Correct & Mean ms & Uploads & Skipped \\
\midrule
UCF101 & SmolVLM & 928/3783 & 350 & 0 & 29.5\% \\
UCF101 & Qwen & 1785/3783 & 348 & 1.00 & 0 \\
UCF101 & Live rule & 1640/3783 & 816 & 1.49 & 29.3\% \\
HMDB51 & SmolVLM & 288/1530 & 373 & 0 & 24.3\% \\
HMDB51 & Qwen & 812/1530 & 371 & 1.00 & 0 \\
HMDB51 & Live rule & 660/1530 & 933 & 1.71 & 24.0\% \\
\bottomrule
\end{tabular}
\end{table}

Table~\ref{tab:video} puts Qwen ahead of the live rule on both lists, by
145 clips on UCF101 and by 152 clips on HMDB51. SmolVLM alone is far
behind both. The quiet skip removes 29.5\% of UCF101 frames and 24.3\% of
HMDB51 frames from the phone model, and the live rule keeps almost the
same skip rate. That saving does not make the live rule cheaper than one
Qwen call. A clip can upload more than once, because an early low
confidence sends the prefix seen so far and a later frame can send again.
The mean is 1.49 uploads on UCF101 and 1.71 on HMDB51, against one upload
for the Qwen branch. The 95th percentile compute is 1.6~s and 1.7~s for
the live rule, against 0.46~s for Qwen on both lists.

A grid of \(\tau_{\mathrm{u}}\in\{0.3,0.4,0.5,0.6,0.7\}\) and
\(\tau_{\mathrm{s}}\in\{0.60,0.65,0.70,0.75,0.80\}\) does not reverse the
order. On UCF101 the live rule moves from 1332/3783 at (0.3, 0.60), with
0.94 uploads and a mean of 648~ms, to 1776/3783 at (0.7, 0.80), with 1.82
uploads and a mean of 913~ms. On HMDB51 the same corners are 487/1530 and
729/1530. The high corner is still short of the single Qwen pass. Raising
both thresholds uploads more often and spends more compute to approach the
model that already saw all three frames.

\subsection{What these measurements do not support}

The 60 images are a phrase-quality check. They do not measure how often
the live gate uploads. The replay measures that count on three stored
frames, not on a phone session and not on the Wi-Fi link. Neither
measurement supports a claim about action start times. They do not compare
schedulers. The class score is also a lower bound on descriptive
usefulness: a phrase such as ``Smiling'' on a drinking image is the wrong
class and may still be a fair description of the visible face. The report
does not treat those misses as correct.
